\chapter*{RESUMEN EJECUTIVO}
\addcontentsline{toc}{chapter}{RESUMEN EJECUTIVO}
\chaptermark{RESUMEN EJECUTIVO}

En la actualidad, la adopción masiva de asistentes de Inteligencia Artificial
Generativa (IAG) como GitHub Copilot, Claude y ChatGPT ha revolucionado los
entornos de desarrollo de software, acelerando la generación inicial de código
fuente. Sin embargo, en dominios altamente regulados y conceptualmente
complejos --como el sector de Seguros y Reaseguros--, el uso irrestricto de
herramientas probabilísticas de codificación suele derivar en una degradación
silenciosa de la arquitectura del software.

Las reglas de negocio críticas (tarificación actuarial, suscripción de
riesgos, coaseguros y bordereaux de reaseguro) requieren alta precisión, bajo
acoplamiento y estricta trazabilidad. La sobre-dependencia de la IAG promueve
la aceptación de código con ``alucinaciones de diseño'', duplicación masiva,
lógica defensiva redundante y desacoplamiento de patrones arquitectónicos,
generando un fenómeno denominado Deuda Técnica Generativa (DTG). El problema
fundamental radica en la falta de un modelo cuantitativo que permita medir y
predecir el impacto real de esta deuda sobre la mantenibilidad y el costo de
vida útil de las aplicaciones empresariales.
